\chapter{Datasets, Segmentation Models and Offline Evaluation}
\label{cap:dataset_modelli}

The segmentation study began with a comparison of FCN, DeepLabV3, and
LR-ASPP on manually annotated lettuce images. Manual segmentation was
time-consuming, so this first comparison was limited to lettuce rather
than repeated for every crop. Its outcome led to the adoption of LR-ASPP.

SAM~3 was then used to generate training labels for lettuce, tomato, and
cucumber images. The resulting multi-crop models differed in their output
classes: plants and path plus background, four classes including covering
material, or path and background alone. These models were evaluated
against the manual lettuce masks. The crop coverage of the training data
therefore differs from that of the reported evaluation set.

This chapter describes the datasets and training procedure, followed by
the offline comparisons that motivate the configuration selected for the
robot. The performance of the trained networks is considered separately
from the quality of the SAM~3 annotations themselves. The architecture results are reported from the available validation
exports; the comparisons of supervision strategies still require checks
on evaluation provenance. Geometric and closed-loop field experiments
are discussed later in the thesis.

\section{Image acquisition and class definition}
\label{sec:c3_acquisition_classes}

\subsection{Image sources and acquisition conditions}
\label{subsec:c3_image_acquisition}

\subsection{Semantic classes and labeling targets}
\label{subsec:c3_semantic_classes}

Three distinct labeling configurations were evaluated for SAM~3:
\begin{itemize}
\item \textbf{Four-class configuration:} Segments background, plants, soil, and covering material. This configuration was also used for the lettuce models during architecture selection.
\item \textbf{Two-foreground-class configuration:} Labels plants and path, with background as the third output class. This produces three total semantic categories.
\item \textbf{Path-only configuration:} Binary segmentation with background and path as the two output labels. For robotic integration, the path mask is the relevant output. Segmenting soil alone does not identify a unique navigation corridor or establish that it is obstacle-free. The geometric methods in Chapter~\ref{cap:geometric_reference} use the mask to construct a local image-space reference.
\end{itemize}

\section{Lettuce, tomato, and cucumber datasets}
\label{sec:c3_datasets}

\subsection{Lettuce dataset}
\label{subsec:c3_lettuce}

The manually annotated lettuce dataset, identified as \texttt{Data},
contains 762 image--mask pairs distributed across eight folders. Its
masks were prepared from the manual annotations stored in
\texttt{annotations.json}. The same manual reference is used for the
initial architecture comparison and for evaluating the models trained
with SAM~3-generated multi-crop labels. Automatically annotated lettuce
images belong to the latter training datasets, not to separate
lettuce-only training configurations.

\subsection{SAM-assisted multi-crop datasets and label configurations}
\label{subsec:c3_multicrop_path}

The automatically supervised models use images spanning all the crops
considered in this work. The change between configurations concerns the
semantic targets to be segmented, not a restriction to a different crop:
plants and path plus background, plants and path and covering material
plus background, or path plus background alone. The exact image lists,
partitions, and dataset export associated with each checkpoint must
nevertheless be retained; common crop coverage does not by itself verify
identical sample membership across training runs.

The binary \texttt{allPlantsSAM} model was trained on
\texttt{Data\_\allowbreak allPlants\allowbreak SAM2}. Its automatic masks
describe only path, with all remaining pixels assigned to background.
This formulation focuses the output on the target region used by the
geometric extractor, rather than on classification of every scene component.
The dataset-directory identifiers for the other label configurations
must be established from their training records, not inferred from the
names of the lettuce evaluation files.

The lettuce evaluation examines these multi-crop models using manually
annotated lettuce images. Multi-crop training does not establish that
those evaluation images are unseen. Possible image overlap, duplicates,
or related acquisitions between training and evaluation have not yet
been checked; conclusions about generalization remain subject to that
verification.

\subsection{Tomato dataset}
\label{subsec:c3_tomato}

The tomato images were labeled using SAM~3 rather than the manual
procedure used for the initial lettuce comparison. They contribute to
the automatically annotated multi-crop training data. LR-ASPP was adopted
after the comparison of the three candidate architectures on lettuce;
no separate FCN, DeepLabV3, and LR-ASPP comparison was performed on tomato.

\subsection{Cucumber dataset}
\label{subsec:c3_cucumber}

SAM~3 was also used to label the cucumber images, which contribute to
the automatically annotated multi-crop training data. As for tomatoes,
the architecture choice followed the comparison on manually annotated
lettuce, rather than a separate comparison on cucumber.

\section{Manual and SAM 3-assisted labeling}
\label{sec:c3_annotation}

\subsection{Manual labeling procedure}
\label{subsec:c3_manual_annotation}

The lettuce images were labeled manually before training the three
candidate architectures. The masks derived from
\texttt{annotations.json} are adopted as the ground truth for evaluating
the trained segmentation models on lettuce. They can also serve as the
reference for a separate direct comparison with SAM~3-generated masks.
Using these annotations as ground truth defines the evaluation reference;
it does not imply that manual annotation is intrinsically error-free.

\subsection{Model-assisted labeling procedure}
\label{subsec:c3_assisted_annotation}

The time required for manual labeling motivated the use of SAM~3 to
prepare the multi-crop datasets. Initial use suggested faster mask
generation, but the total annotation effort has not been measured in a
systematic comparison that includes inspection and correction.

The SAM~3-assisted labeling procedure covers lettuce, tomato, and
cucumber images for the multi-crop training configurations. In this
workflow, automatic segmentation is used to prepare labels offline; it
is distinct from the segmentation model used during navigation.

For the lettuce images annotated with two foreground classes, the reported
text prompts are \texttt{plants} and \texttt{central farm path}. The
four-class labeling additionally uses \texttt{plastic mulch film} for
covering material. Background accounts for the remaining label in each
configuration. These reported lettuce prompts do not establish the exact
prompts used for every crop in the multi-crop datasets. The binary
configuration retains only the path annotation; crop-specific prompts
and mask-composition settings for each configuration must be documented
from the corresponding labeling runs.

These text descriptions define annotation targets, but do not specify
how overlapping masks are resolved. The available
\texttt{sam3\_\allowbreak autolabel\_\allowbreak 3class.py} is configured
for tomatoes with the prompts \texttt{tomato plant} and \texttt{soil};
it gives plant pixels precedence over soil pixels when composing the
mask. Its settings should not be attributed to the lettuce labeling
runs without checking their implementation. In particular, an annotation
of visible soil and a prompt requesting the central farm path can target
different regions. Agreement with the manual soil class must therefore
be assessed rather than assumed.

Initial visual inspection suggested good results on images considered
relatively simple, although this impression still requires systematic
evaluation. Inspection of the lettuce masks also revealed cases in which the
automatic procedure identified only some of the plants visible in the
image. Other masks were less precise than their manual counterparts.
Errors were observed particularly in images containing small lettuce
plants, while distant plants were a recurring difficulty among the
images with labeling errors. These are qualitative observations from
the inspected images, not yet measurements of error frequency or evidence
that the same behavior occurs to the same extent in the other crops.

\section{Labeling quality and dataset partitioning}
\label{sec:c3_quality_splits}

\subsection{Label inspection and revision criteria}
\label{subsec:c3_annotation_quality}

The planned quantitative assessment evaluates the quality of the
SAM~3-generated lettuce masks against the manual ground-truth masks.
The intended measure is the intersection over union for the lettuce
class, comparing the automatically labeled pixels with the corresponding
manual reference. This comparison is intended to quantify the quality
of SAM~3 labeling on the evaluated lettuce images, not its performance
across all crops or applications. The exact evaluation set and aggregation
procedure have not yet been fixed, and no numerical labeling-quality
results are reported here.

This assessment must be kept separate from the comparison of the trained
segmentation models. Agreement between two sets of labels does not
directly measure the performance of a network trained on either set.
The additional models trained on automatic multi-crop labels address a
separate question: how the resulting segmentation networks perform
against a common manual reference. Their performance also depends on
the training images, target classes, model configuration, and optimization.
The difference between two networks cannot therefore be attributed
exclusively to annotation quality when other factors differ.

\subsection{Training, validation, and test sets}
\label{subsec:c3_dataset_splits}

The reported training runs use an 80\% training, 10\% validation, and
10\% test partition with seed 42. The available \texttt{train\_model.py}
implements these proportions by assigning the integer parts of the
training and validation counts and leaving the remaining samples for
test. Applied to 762 pairs, this rule gives 609 training, 76 validation,
and 77 test samples. These counts describe the intended manual-dataset
partition; the saved indices and image lists must establish the actual
membership of each evaluated split.

The test set is not used for gradient-based parameter updates. Checkpoint
selection uses validation rather than test performance. A common seed
alone does not ensure that differently sized or ordered datasets yield
identical image partitions. Corresponding image identities must therefore
be checked when comparing the manually supervised lettuce model with the
automatically supervised multi-crop models, including whether a lettuce
evaluation image appears in an automatic model's training set.

The reported operational evaluation selects 42 images from the manual
test split belonging to \texttt{insalata\_grande1},
\texttt{insalata\_grande2}, and \texttt{insalata\_media1}. These folders
represent the row conditions selected for the robot application. The
42 images are a subset of the reported 77 test images, not an additional
independent test set. Performance on this subset describes those selected
conditions, rather than the full range of lettuce images.

\section{FCN, DeepLabV3, and LR-ASPP: initialization and fine-tuning}
\label{sec:c3_segmentation_models}

\subsection{Selection of the TorchVision implementations}
\label{subsec:c3_model_selection}

FCN, DeepLabV3, and LR-ASPP were chosen from the implementations available
in TorchVision. All three were trained on the manually annotated lettuce
dataset. Their selection was practical, rather than the outcome of a
comparison with every available architecture. Accuracy and inference
time were then used to choose among these candidates, as discussed in
Section~\ref{sec:c3_architecture_comparison}.

\subsection{Pretrained weights and adaptation of the output classes}
\label{subsec:c3_pretrained_weights}

The initialization and training strategy described below follows the
available implementation in \texttt{models.py} and
\texttt{train\_model.py}. The current training script is configured for
the tomato dataset; its correspondence with the earlier lettuce runs
remains to be confirmed before these settings are associated with their
experimental results.

The implementation separates each network into a backbone, which extracts
image features, and a segmentation head, which produces class scores at
each output location. Transfer learning is implemented by loading
pretrained parameters and adapting the output layers to the target
labeling task, rather than initializing every network parameter from
scratch. The model construction function accepts the number of output
classes as an argument; the available training script explicitly sets
this number to three.

For FCN and DeepLabV3, the selected implementations use a ResNet-50
backbone. They are constructed with the default pretrained segmentation
weights specified by \texttt{FCN\_ResNet50\_Weights} and
\texttt{DeepLabV3\_ResNet50\_Weights}, respectively. After these weights
are loaded, the final convolution of the main classifier is replaced by
a new $1\times1$ convolution with one output channel per target class.
The remaining pretrained layers are retained. Both models also include
an auxiliary classifier, whose final convolution is replaced in the same
way. Thus, the output layers associated with the original classes are
newly initialized, while the preceding layers retain their pretrained
parameters.

LR-ASPP is initialized differently. The implementation uses a
MobileNetV3-Large backbone and requests its default pretrained weights
through the \texttt{weights\_\allowbreak backbone} argument. However, the
\texttt{weights} argument of the complete segmentation model is set to
\texttt{None}. Consequently, the backbone is pretrained, whereas the
LR-ASPP segmentation head is newly initialized for the requested number
of classes. The code does not load a pretrained LR-ASPP segmentation
checkpoint and then replace only its final layers.

This distinction is relevant to the interpretation of the model
comparison. The three configurations differ not only in architecture,
but also in the extent of their pretrained initialization: FCN and
DeepLabV3 retain pretrained segmentation-head layers, whereas LR-ASPP
starts with a new segmentation head. Their results therefore describe
the implemented configurations, not an isolated comparison of
architectures under identical initialization conditions.

\subsection{Backbone freezing and fine-tuning}
\label{subsec:c3_fine_tuning}

Training is organized into two stages. Before the first epoch, the
backbone parameters are excluded from gradient-based updates by setting
their \texttt{requires\_grad} attribute to \texttt{False}. The complete
main classifier remains trainable, as does the auxiliary classifier when
present. This first stage therefore updates the segmentation heads,
not just the newly replaced output convolutions of FCN and DeepLabV3.

For the first ten epochs, AdamW is constructed using only the trainable
parameters, with an initial learning rate of $10^{-3}$ and weight decay
of $10^{-4}$. A cosine-annealing scheduler is stepped after each epoch,
with a minimum learning rate of $10^{-6}$. Its initial horizon is set to
the configured total of 200 epochs, although this scheduler is replaced
at the transition to the second stage.

At the end of epoch ten, the script enables gradients for all backbone
parameters. Fine-tuning of the backbone therefore begins at epoch eleven,
together with continued training of the segmentation heads. A new AdamW
optimizer is constructed with separate parameter groups: the backbone
starts this stage at a learning rate of $10^{-5}$, while the main and,
where present, auxiliary classifiers start at $10^{-4}$. Weight decay
remains $10^{-4}$. The backbone learning rate is thus lower than that of
the heads, allowing the pretrained features to be adapted with smaller
learning-rate settings during this stage.

The optimizer is recreated rather than extended in place, so its
accumulated state from the first stage is not transferred to the second.
The model parameters themselves are retained at the transition. A new
cosine-annealing scheduler is also created for the remaining 190 epochs,
again with a minimum learning rate of $10^{-6}$. The script configures
200 epochs in total; the early-stopping block is commented out and does
not terminate training in the active implementation.

Here, backbone freezing refers specifically to parameter updates. The
training loop calls \texttt{model.train()} without separately placing the
backbone in evaluation mode. Its batch-normalization running statistics
therefore continue to update during training, including the first stage.
The backbone is not completely fixed in the sense of preserving all of
its internal state. Validation is instead performed with
\texttt{model.eval()} and gradient computation disabled.

This two-stage procedure first adapts the segmentation heads and then
jointly fine-tunes the backbone and heads. The loss function, data
transformations, and checkpoint-selection criterion are addressed in
Section~\ref{sec:c3_training}, separately from the initialization and
parameter-update schedule described here.

\section{Training configuration and checkpoint selection}
\label{sec:c3_training}

\subsection{Preprocessing and data augmentation}
\label{subsec:c3_preprocessing}

\subsection{Loss function and optimization settings}
\label{subsec:c3_loss_optimization}

\subsection{Validation and checkpoint selection}
\label{subsec:c3_checkpoint_selection}

In the available training implementation, validation is performed without
gradient updates after each epoch. A checkpoint is saved as
\texttt{best.pth} whenever the validation loss is lower than all previous
values; \texttt{last.pth} stores the latest training state. Selection
therefore uses minimum validation loss, not maximum test IoU or simply
the last epoch. The script records the epoch, training and validation
losses, split proportions, and seed with the checkpoint, and saves split
indices separately.

This describes the supplied training script, currently configured for
tomatoes. The corresponding records must confirm the criterion and
settings used for the manual lettuce checkpoint and each SAM-supervised
multi-crop checkpoint evaluated on lettuce.
The differing checkpoint epochs in the comparison files do not, on their
own, establish either a common training duration or a causal explanation
for differences in segmentation accuracy.

\section{Offline architecture comparison on manually annotated lettuce}
\label{sec:c3_architecture_comparison}

This comparison addresses the architecture choice before changing the
annotation procedure or crop coverage. FCN, DeepLabV3, and LR-ASPP were
trained and evaluated using the manually annotated lettuce dataset.
The implemented initializations differ as described in
Section~\ref{subsec:c3_pretrained_weights}; the comparison concerns these
trained configurations rather than architecture alone under identical
pretraining conditions.

\subsection{Evaluation protocol and inference-time measurement}
\label{subsec:c3_architecture_protocol}

The comparison uses the three CSV records stored in
\texttt{training insalata}: the FCN, DeepLabV3, and LR-ASPP
\texttt{*\_\allowbreak validation\_\allowbreak metrics.csv} exports.
Each contains aggregate metrics, the checkpoint epoch, and IoU values
for background, plants, soil, and covering material. The reported mIoU
agrees with the arithmetic mean of these four class IoUs, including
background. The records are presented here as validation exports,
not as results of an independent final test.

The CSV files do not identify the evaluated images or timing hardware.
The supplied \texttt{validation\_model.py} selects validation indices,
aggregates a confusion matrix, and measures the model forward pass and
\texttt{argmax} at batch size one, with CUDA synchronization before and
after the timed block when a GPU is available. It contains no explicit
warm-up stage. However, that script is currently configured for tomato
images and three classes, whereas these exports contain four classes.
The original lettuce evaluation configuration must therefore confirm the
image lists, hardware, resolution, and timing procedure before the
recorded speed ratios are treated as a fully controlled benchmark.

\subsection{Accuracy and inference-time results}
\label{subsec:c3_architecture_results}

Table~\ref{tab:c3_architecture_results} reports the exported metrics.
Inference times have been converted from seconds to milliseconds.
The FPS values equal the reciprocal of the mean recorded time; they
are not measurements of the complete robot processing or control rate.
Checkpoint epochs identify the evaluated weights, not training duration.

\begin{table}[htbp]
    \centering
    \small
    \caption{Architecture comparison from the manual lettuce validation
    exports. Timing values are those recorded in the CSV files, not
    end-to-end robot latency.}
    \label{tab:c3_architecture_results}
    \begin{tabular}{@{}l r r r r r@{}}
        \toprule
        Model & Epoch & mIoU & Path IoU &
        \shortstack{Mean time\\(ms)} & FPS \\
        \midrule
        FCN & 38 & 0.9232 & 0.9621 & 58.24 & 17.17 \\
        DeepLabV3 & 43 & 0.9228 & 0.9651 & 86.87 & 11.51 \\
        LR-ASPP & 194 & 0.9108 & 0.9546 & 10.41 & 96.07 \\
        \bottomrule
    \end{tabular}
\end{table}

FCN and DeepLabV3 achieve mIoU values of 0.9232 and 0.9228, respectively,
compared with 0.9108 for LR-ASPP. Using the unrounded CSV values, the
LR-ASPP deficit is 1.24 percentage points relative to FCN and 1.19
relative to DeepLabV3. Its mean recorded time is 10.41~ms, compared
with 58.24~ms and 86.87~ms. These timings correspond to a 5.60-fold
reduction relative to FCN and an 8.35-fold reduction relative to
DeepLabV3.

Table~\ref{tab:c3_architecture_class_iou} provides the class-level
comparison. LR-ASPP has lower IoU than both alternatives in each class,
although its soil/path IoU remains 0.9546. The data support a trade-off
between accuracy and latency, not a claim of equal segmentation accuracy.

\begin{table}[htbp]
    \centering
    \small
    \caption{Per-class IoU from the same manual lettuce validation
    exports. Plants, soil, and covering material correspond to
    \texttt{pianta}, \texttt{suolo}, and \texttt{copertura} in these files.}
    \label{tab:c3_architecture_class_iou}
    \begin{tabular}{@{}l r r r r@{}}
        \toprule
        Model & Background & Plants & Soil/path & Covering material \\
        \midrule
        FCN & 0.9464 & 0.9280 & 0.9621 & 0.8565 \\
        DeepLabV3 & 0.9464 & 0.9277 & 0.9651 & 0.8519 \\
        LR-ASPP & 0.9394 & 0.9120 & 0.9546 & 0.8372 \\
        \bottomrule
    \end{tabular}
\end{table}

\subsection{Motivation for adopting LR-ASPP}
\label{subsec:c3_lraspp_selection}

LR-ASPP was adopted because the substantial reduction in recorded
inference time justified accepting an mIoU approximately 1.2 percentage
points below the other two configurations. It was not the most accurate
model, but offered the practical accuracy--latency compromise sought for
robot perception. FCN and DeepLabV3 had higher IoU but substantially
longer recorded inference times. Subsequent SAM-supervised multi-crop
training therefore used LR-ASPP while varying the output label set.

This choice is based on the available offline comparison. The records
do not establish a hardware-independent speed ranking or demonstrate
that the integrated robot meets its control timing requirements.
Those requirements must be assessed using the deployed model and the
complete perception and control pipeline.

\section{Offline comparison of supervision strategies}
\label{sec:c3_supervision_comparison}

\subsection{Training configurations and evaluation targets}
\label{subsec:c3_supervision_configurations}

Table~\ref{tab:c3_supervision_configurations} summarizes the four training
configurations considered in this comparison. This study is separate
from the initial comparison of the three candidate architectures.
The manual configuration is trained on lettuce only, whereas all three
SAM-supervised configurations use multi-crop data and differ in their
output label sets. The model names identify supervision and targets;
they do not mean that SAM~3 itself is executed during robot navigation.
Each dataset export, image partition, architecture, and checkpoint must
be retained in the experiment records.

\begin{table}[htbp]
    \centering
    \small
    \caption{Training scope and target labels for the comparison of manual
    and SAM~3-generated supervision. Class counts include background;
    evaluation uses the manual lettuce reference.}
    \label{tab:c3_supervision_configurations}
    \begin{tabularx}{\textwidth}{@{}X X X@{}}
        \toprule
        Configuration & Training scope & Output labels \\
        \midrule
        Manual & Lettuce only (\texttt{Data}) &
        Background, plants, path, covering material \\
        SAM-generated, two foreground classes &
        Multi-crop & Background, plants, path \\
        SAM-generated, four classes &
        Multi-crop & Background, plants, path, covering material \\
        SAM-generated, binary (\texttt{allPlantsSAM}) &
        Multi-crop & Background, path \\
        \bottomrule
    \end{tabularx}
\end{table}

Five evaluations on manually annotated lettuce images are organized
around these configurations:
\begin{enumerate}
    \item the manual and automatically supervised four-class models,
    evaluated on all four categories over the full evaluation split;
    \item the manual model and the automatically supervised model with
    two foreground classes, evaluated on plants and path;
    \item the manual and SAM-supervised binary models, evaluated only on path;
    \item the same path-only comparison on the selected operational subset;
    \item the four-class comparison on that operational subset.
\end{enumerate}
The second evaluation uses a separate automatically supervised
checkpoint, not merely a different average of the four-class model's
metrics. The subset evaluations reuse their respective trained models;
changing the evaluation images does not constitute retraining.

Manual masks supply the ground truth for all comparisons. For path-only
scoring, the target is the manual soil class and the predicted positive
region is the corresponding path output of each model. This definition
allows a four-class and a binary model to be evaluated on the same
positive class. The treatment of other manual classes, including covering
material, must be explicit in the evaluation implementation; they must
not be silently omitted when defining the pixels used for scoring.

\subsection{Metrics and interpretation of the comparison}
\label{subsec:c3_supervision_metrics}

For a class $c$, intersection over union is defined as
\begin{equation}
    \operatorname{IoU}_c=
    \frac{\operatorname{TP}_c}
         {\operatorname{TP}_c+\operatorname{FP}_c+\operatorname{FN}_c},
    \label{eq:c3_iou}
\end{equation}
where true positives, false positives, and false negatives are defined
against the manual reference. The available \texttt{metrics.py} computes
class IoUs from a confusion matrix. The aggregation used by each exported
comparison must nevertheless be confirmed: computing IoU from pooled
pixel counts and averaging image-wise IoUs are not interchangeable.

The four-class mIoU averages background, plants, path, and covering
material. The two-foreground-class comparison instead averages only
plants and path; the path-only comparison reports a single class IoU.
These averages concern different targets and are not used as directly
comparable overall rankings. Path IoU is the common task-specific measure,
provided the evaluation images and treatment of ground-truth pixels
are aligned.

Before completing the numerical tables in this section as full-test
results, the evaluation provenance must be
resolved: the currently available four-class and two-foreground-class
CSV files have \texttt{val} in their names and contain values different
from the reported full-test summary. Their filenames alone do not prove
which split was used. Saved image lists, scripts, and logs are needed to
reconcile those exports with the reported 77-image test protocol.
Inference-time comparisons likewise require consistent hardware,
preprocessing, and timing boundaries.

\subsection{Results on the full lettuce evaluation set}
\label{subsec:c3_supervision_full_results}

\emph{Results to complete after provenance checks:} Present the four-class,
two-foreground-class, and path-only comparisons on the documented lettuce
evaluation set. Report per-class IoU and state which classes each average
includes. Add a path-focused summary, identifying the checkpoint and split
for each configuration. Keep validation exports separate from verified
test results rather than combining their values.

\subsection{Results on the selected operational subset}
\label{subsec:c3_supervision_operational_results}

\emph{Results to complete after image-list verification:} Present the
path-only and four-class comparisons on the reported 42-image operational
subset, alongside the manual reference. Identify the included folders
and the checkpoints reused from the full-set evaluation. Interpret
changes as differences between evaluated image groups, not improvements
caused by retraining. The subset is not an independent additional test.

\subsection{Discussion and limitations of the offline comparison}
\label{subsec:c3_supervision_discussion}

Training crop coverage, supervision source, and output classes must be
kept distinct when interpreting the results. The manual model and the
SAM-supervised models differ in more than annotation source: their
training images also differ in crop coverage. Their performance gap
cannot therefore be attributed solely to annotation quality. Likewise,
a path-only model is not intrinsically more accurate because its label
set matches the downstream task.

Claims about generalization require checking image overlap and related
acquisitions. Performance on the selected operational subset does not
establish the same accuracy across the full dataset. These offline
comparisons do not establish geometric-reference accuracy, obstacle-free
passage, or closed-loop navigation performance.

\section{Segmentation configuration selected for robotic integration}
\label{sec:c3_deployment_selection}

The two selection decisions are complementary. The initial manual
lettuce comparison motivates the architecture, LR-ASPP. The later
multi-crop comparison concerns supervision and target classes, leading
to the binary path configuration referred to as \texttt{allPlantsSAM}.
The latter decision does not replace the architecture comparison.

The binary configuration was selected because its target is directly
aligned with path extraction and its observed behavior on the selected
operational images supported that choice. The quantitative evidence and
limitations are to be completed in
Section~\ref{sec:c3_supervision_comparison}. This is not a claim that the
binary configuration outperforms the manually supervised model or that
performance on the selected rows holds across the entire dataset.

\emph{Deployment record to complete:} Identify the exact trained checkpoint
and exported model used by the robot, input resolution, preprocessing,
and path-label encoding. Relate its output to the geometric interface in
Chapter~\ref{cap:geometric_reference}. Report quantitative comparisons
only after their evaluation provenance has been established.

If test results contributed to selecting the deployed configuration,
disclose that role: those images then served model selection rather than
an entirely independent final assessment. Later geometric and field
experiments must separately assess reference quality and the behavior
of the integrated robot; segmentation IoU alone does not demonstrate
successful row following.

\section{ONNX export and inference configuration}
\label{sec:c3_export_inference}

\subsection{Model export}
\label{subsec:c3_onnx_export}

\subsection{Inference inputs and outputs}
\label{subsec:c3_inference_interface}
\mbox{}\par